\documentclass[a4paper,12pt]{article}

\usepackage[T2A]{fontenc}
\usepackage[utf8]{inputenc}
\usepackage[russian]{babel}

\usepackage{amsmath, amssymb, amsthm, mathtools}
\usepackage{helvet}
\renewcommand{\familydefault}{\sfdefault}
\usepackage{icomma}
\usepackage{geometry}
\usepackage{microtype}
\usepackage{tikz}
\usetikzlibrary{calc,arrows.meta}
\usepackage{tkz-euclide}
\usepackage{pgfplots}
\pgfplotsset{compat=1.18}
\usepackage{tabularx, booktabs, longtable, array}
\usepackage{enumitem}
\usepackage{hyperref}
\usepackage{parskip}
\usepackage{fancyhdr}
\usepackage{setspace}
\usepackage{xcolor}

\geometry{left=20mm,right=20mm,top=20mm,bottom=22mm}
\onehalfspacing

\definecolor{accent}{HTML}{008080}
\definecolor{darkaccent}{HTML}{004D4D}
\definecolor{textgray}{HTML}{333333}
\definecolor{softgray}{HTML}{F2F5F5}

\hypersetup{
    colorlinks=true,
    linkcolor=darkaccent,
    urlcolor=darkaccent
}

\pagestyle{fancy}
\fancyhf{}
\lhead{\textcolor{textgray}{Логарифмические неравенства}}
\rhead{\textcolor{textgray}{ЕГЭ}}
\cfoot{\thepage}
\renewcommand{\headrulewidth}{0.4pt}

\setlist[itemize]{leftmargin=1.2cm,itemsep=0.15em,topsep=0.25em}
\setlist[enumerate]{leftmargin=1.2cm,itemsep=0.25em,topsep=0.25em}

\providecommand{\tg}{\mathop{\mathrm{tg}}\nolimits}
\providecommand{\ctg}{\mathop{\mathrm{ctg}}\nolimits}
\providecommand{\arctg}{\mathop{\mathrm{arctg}}\nolimits}

\newcommand{\topic}{Логарифмические неравенства: ОДЗ, основания, модули и метод интервалов}
\newcommand{\important}[1]{\textcolor{darkaccent}{\textbf{#1}}}

\begin{document}

\begin{titlepage}
\thispagestyle{empty}
\begin{center}
\vspace*{3.2cm}

{\fontsize{30}{36}\selectfont\textbf{Чек-лист}}

\vspace{0.8cm}

{\Large\textcolor{darkaccent}{\textbf{\topic}}}

\vspace{2.4cm}

{\large Лёвин Артём Александрович эксклюзивно для Николь}

\vspace{0.7cm}

{\large ЕГЭ}

\vspace{0.7cm}

{\large \today}

\vfill

\begin{tikzpicture}[remember picture,overlay]
\draw[
    accent,
    line width=1.2pt
]
([xshift=1.8cm,yshift=2.1cm]current page.south west)
.. controls
([xshift=5.2cm,yshift=3.1cm]current page.south west)
and
([xshift=-5.2cm,yshift=1.1cm]current page.south east)
..
([xshift=-1.8cm,yshift=2.1cm]current page.south east);
\draw[
    darkaccent,
    line width=0.6pt
]
([xshift=3.0cm,yshift=1.55cm]current page.south west)
.. controls
([xshift=7.0cm,yshift=2.25cm]current page.south west)
and
([xshift=-7.0cm,yshift=0.95cm]current page.south east)
..
([xshift=-3.0cm,yshift=1.55cm]current page.south east);
\end{tikzpicture}

\end{center}
\end{titlepage}

\tableofcontents
\newpage

\section{Главная идея темы}

Логарифмические неравенства решаются не только ``отбрасыванием логарифмов''. Перед этим нужно проверить основание логарифма, область допустимых значений и структуру выражения под логарифмом.

\subsection*{Базовая схема}

\begin{enumerate}
    \item Запишите \important{ОДЗ}: все выражения под логарифмами должны быть строго положительными.
    \item Приведите числа к логарифмам с тем же основанием, если это удобно:
    \[
        c=\log_a a^c,\qquad a>0,\quad a\ne 1.
    \]
    \item Сравните основание с единицей:
    \[
    \begin{cases}
    a>1 & \Rightarrow \text{знак неравенства сохраняется},\\
    0<a<1 & \Rightarrow \text{знак неравенства меняется}.
    \end{cases}
    \]
    \item После снятия логарифмов решайте полученное алгебраическое неравенство.
    \item Если появляется дробь, приводите её к виду:
    \[
        \frac{\text{многочлен}}{\text{многочлен}}\ \vee\ 0,
    \]
    где справа стоит ноль, а слева одна дробь.
    \item Решайте рациональное неравенство методом интервалов.
    \item В конце обязательно сверяйте ответ с ОДЗ.
\end{enumerate}

\section{Основание логарифма}

\subsection{Когда знак сохраняется}

Если \(a>1\), то функция \(y=\log_a x\) возрастает. Поэтому:
\[
    \log_a f(x) \ge \log_a g(x)
    \quad \Longleftrightarrow \quad
    \begin{cases}
    f(x)>0,\\
    g(x)>0,\\
    f(x)\ge g(x).
    \end{cases}
\]

Пример:
\[
    \log_4 A \ge \frac12.
\]
Так как
\[
    \frac12=\log_4 2,
\]
получаем:
\[
    \log_4 A \ge \log_4 2
    \quad \Longleftrightarrow \quad
    A\ge 2,
    \qquad A>0.
\]
При \(A\ge 2\) условие \(A>0\) выполняется автоматически.

\subsection{Когда знак меняется}

Если \(0<a<1\), то функция \(y=\log_a x\) убывает. Поэтому:
\[
    \log_a f(x) \ge \log_a g(x)
    \quad \Longleftrightarrow \quad
    \begin{cases}
    f(x)>0,\\
    g(x)>0,\\
    f(x)\le g(x).
    \end{cases}
\]

Пример:
\[
    \log_{\frac12} f(x) \le \log_{\frac12} g(x)
    \quad \Longleftrightarrow \quad
    \begin{cases}
    f(x)>0,\\
    g(x)>0,\\
    f(x)\ge g(x).
    \end{cases}
\]

\subsection{Как избавиться от дробного основания}

Если основание меньше единицы, иногда удобно переписать его через отрицательную степень:
\[
    \frac12=2^{-1},\qquad
    \frac34=\left(\frac43\right)^{-1}.
\]

Полезная формула:
\[
    \log_{a^{-1}} b=-\log_a b.
\]

Например:
\[
    \log_{\frac12} u=\log_{2^{-1}}u=-\log_2 u.
\]

Если всё неравенство умножается на \(-1\), знак обязательно меняется:
\[
    -\log_2 f(x)\ge -\log_2 g(x)
    \quad \Longleftrightarrow \quad
    \log_2 f(x)\le \log_2 g(x).
\]

\section{ОДЗ в логарифмических неравенствах}

\subsection{Главное правило}

Выражение под логарифмом должно быть строго положительным:
\[
    \log_a F(x) \quad \Rightarrow \quad F(x)>0.
\]

Если под логарифмом стоит дробь:
\[
    \log_a \frac{P(x)}{Q(x)},
\]
то нужно требовать:
\[
    \frac{P(x)}{Q(x)}>0.
\]
Отдельно писать \(Q(x)\ne 0\) можно, но обычно это уже учитывается при решении дробного неравенства.

\subsection{ОДЗ может упростить задачу}

Иногда ОДЗ не просто ограничивает ответ, а заранее показывает, как раскрывать модуль.

Например:
\[
    \log_4 \frac{4x-5}{|x-1|}>\frac12.
\]
ОДЗ:
\[
    \frac{4x-5}{|x-1|}>0.
\]
Так как \(|x-1|>0\) при \(x\ne 1\), знак дроби определяется числителем:
\[
    4x-5>0
    \quad \Longleftrightarrow \quad
    x>\frac54.
\]
На этом множестве \(x>1\), значит:
\[
    |x-1|=x-1.
\]
Поэтому модуль раскрывается только одним способом.

\section{Метод интервалов}

\subsection{Стандартный вид}

Рациональное неравенство нужно привести к виду:
\[
    \frac{P(x)}{Q(x)}\ge 0,
    \qquad
    \frac{P(x)}{Q(x)}>0,
    \qquad
    \frac{P(x)}{Q(x)}\le 0,
    \qquad
    \frac{P(x)}{Q(x)}<0.
\]

Нельзя решать неравенство в виде:
\[
    \frac{P(x)}{Q(x)}\ge 2.
\]
Сначала переносим всё влево:
\[
    \frac{P(x)}{Q(x)}-2\ge 0,
\]
затем приводим к общему знаменателю.

\subsection{Алгоритм}

\begin{enumerate}
    \item Найдите нули числителя.
    \item Найдите нули знаменателя.
    \item Нули числителя отмечайте закрашенными точками, если неравенство нестрогое.
    \item Нули знаменателя всегда выколоты.
    \item Расставьте точки на числовой прямой в правильном порядке.
    \item Найдите знак дроби на одном промежутке подстановкой удобного числа.
    \item Расставьте остальные знаки с учётом кратности корней.
    \item Выберите промежутки нужного знака.
\end{enumerate}

\subsection{Пример с числовой прямой}

Рассмотрим неравенство:
\[
    \log_4\frac{x^2-4}{x+3}\ge \frac12.
\]
Так как \(\frac12=\log_4 2\), а \(4>1\), получаем:
\[
    \frac{x^2-4}{x+3}\ge 2.
\]
Приведём к стандартному виду:
\[
    \frac{x^2-4}{x+3}-2\ge 0,
\]
\[
    \frac{x^2-4-2(x+3)}{x+3}\ge 0,
\]
\[
    \frac{x^2-2x-10}{x+3}\ge 0.
\]
Корни числителя:
\[
    x=1-\sqrt{11},
    \qquad
    x=1+\sqrt{11}.
\]
Корень знаменателя:
\[
    x=-3.
\]
Важно сравнить:
\[
    -3<1-\sqrt{11}<1+\sqrt{11}.
\]

\begin{center}
\begin{tikzpicture}[x=1cm,y=1cm,>=Stealth]
    \draw[->,line width=0.8pt] (-4.4,0) -- (5.4,0);
    \node[below] at (-3,0) {\(-3\)};
    \node[below] at (-2.25,0) {\(1-\sqrt{11}\)};
    \node[below] at (4.25,0) {\(1+\sqrt{11}\)};

    \draw (-3,0.08)--(-3,-0.08);
    \draw (-2.25,0.08)--(-2.25,-0.08);
    \draw (4.25,0.08)--(4.25,-0.08);

    \draw[accent,line width=1.8pt] (-3,0.35)--(-2.25,0.35);
    \draw[accent,line width=1.8pt] (4.25,0.35)--(5.1,0.35);

    \draw[fill=white,line width=0.9pt] (-3,0.35) circle (2.2pt);
    \draw[fill=accent,draw=accent] (-2.25,0.35) circle (2.2pt);
    \draw[fill=accent,draw=accent] (4.25,0.35) circle (2.2pt);

    \node[above] at (-3.6,0.1) {\(-\)};
    \node[above] at (-2.63,0.1) {\(+\)};
    \node[above] at (1,0.1) {\(-\)};
    \node[above] at (4.75,0.1) {\(+\)};
\end{tikzpicture}
\end{center}

Ответ:
\[
    x\in \left(-3;\,1-\sqrt{11}\right]\cup
    \left[1+\sqrt{11};\,+\infty\right).
\]

\section{Модули в логарифмических неравенствах}

\subsection{Модуль как аргумент логарифма}

Если стоит:
\[
    \log_a |F(x)|,
\]
то ОДЗ:
\[
    |F(x)|>0
    \quad \Longleftrightarrow \quad
    F(x)\ne 0.
\]
При этом нужно не забывать про знаменатели внутри \(F(x)\).

Например:
\[
    \log_3\left|\frac{x-2}{x+1}\right|
\]
имеет ограничения:
\[
    x\ne 2,\qquad x\ne -1.
\]

\subsection{Равносильные переходы для модуля}

Если \(g(x)>0\), то:
\[
    |f(x)|>g(x)
    \quad \Longleftrightarrow \quad
    \left[
    \begin{array}{l}
    f(x)>g(x),\\
    f(x)<-g(x).
    \end{array}
    \right.
\]

Если \(g(x)>0\), то:
\[
    |f(x)|<g(x)
    \quad \Longleftrightarrow \quad
    \begin{cases}
    f(x)<g(x),\\
    f(x)>-g(x).
    \end{cases}
\]

Для нестрогих знаков правило сохраняется:
\[
    |f(x)|\ge g(x)
    \quad \Longleftrightarrow \quad
    \left[
    \begin{array}{l}
    f(x)\ge g(x),\\
    f(x)\le -g(x).
    \end{array}
    \right.
\]
\[
    |f(x)|\le g(x)
    \quad \Longleftrightarrow \quad
    \begin{cases}
    f(x)\le g(x),\\
    f(x)\ge -g(x).
    \end{cases}
\]

\subsection{Система и совокупность}

\important{Система} означает одновременное выполнение условий:
\[
    \begin{cases}
    A,\\
    B.
    \end{cases}
\]
В ответ берём пересечение множеств.

\important{Совокупность} означает, что достаточно одного условия:
\[
    \left[
    \begin{array}{l}
    A,\\
    B.
    \end{array}
    \right.
\]
В ответ берём объединение множеств.

Это особенно важно в задачах с модулем:
\[
    |f(x)|>g(x)
\]
даёт совокупность, а
\[
    |f(x)|<g(x)
\]
даёт систему.

\section{Раскрытие модуля по определению}

\[
    |u|=
    \begin{cases}
    u, & u\ge 0,\\
    -u, & u<0.
    \end{cases}
\]

Например:
\[
    |x-1|=
    \begin{cases}
    x-1, & x\ge 1,\\
    1-x, & x<1.
    \end{cases}
\]

Если из ОДЗ уже известно, что \(x>\frac54\), то автоматически \(x>1\), значит:
\[
    |x-1|=x-1.
\]
Это экономит время и защищает от лишних ветвей решения.

\section{Короткие образцы решений}

\subsection*{Образец 1. Основание больше единицы}

\[
    \log_4\frac{x^2-4}{x+3}\ge \frac12.
\]
\[
    \frac12=\log_4 2,
\]
\[
    \frac{x^2-4}{x+3}\ge 2.
\]
\[
    \frac{x^2-4-2(x+3)}{x+3}\ge 0,
\]
\[
    \frac{x^2-2x-10}{x+3}\ge 0.
\]
Ответ:
\[
    x\in \left(-3;\,1-\sqrt{11}\right]\cup
    \left[1+\sqrt{11};+\infty\right).
\]

\subsection*{Образец 2. Основание меньше единицы}

\[
    \log_{\frac12}(x^2-4x+3)\le \log_{\frac12}(x-1).
\]
Так как \(0<\frac12<1\), знак меняется:
\[
    \begin{cases}
    x^2-4x+3>0,\\
    x-1>0,\\
    x^2-4x+3\ge x-1.
    \end{cases}
\]
\[
    \begin{cases}
    (x-1)(x-3)>0,\\
    x>1,\\
    x^2-5x+4\ge 0.
    \end{cases}
\]
\[
    \begin{cases}
    x>3,\\
    (x-1)(x-4)\ge 0.
    \end{cases}
\]
Ответ:
\[
    x\in [4;+\infty).
\]

\subsection*{Образец 3. Модуль под логарифмом}

\[
    \log_3\left|\frac{x-2}{x+1}\right|>1.
\]
\[
    1=\log_3 3,
\]
поэтому:
\[
    \left|\frac{x-2}{x+1}\right|>3.
\]
Получаем совокупность:
\[
    \left[
    \begin{array}{l}
    \dfrac{x-2}{x+1}>3,\\[0.5em]
    \dfrac{x-2}{x+1}<-3.
    \end{array}
    \right.
\]
Первое неравенство:
\[
    \frac{x-2-3(x+1)}{x+1}>0,
    \qquad
    \frac{-2x-5}{x+1}>0.
\]
Второе неравенство:
\[
    \frac{x-2+3(x+1)}{x+1}<0,
    \qquad
    \frac{4x+1}{x+1}<0.
\]
Ответ:
\[
    x\in \left(-\frac52;-1\right)\cup
    \left(-1;-\frac14\right).
\]

\section{Типичные ошибки}

\begin{longtable}{>{\raggedright\arraybackslash}p{0.42\linewidth}>{\raggedright\arraybackslash}p{0.50\linewidth}}
\toprule
\textbf{Ошибка} & \textbf{Как правильно} \\
\midrule
Писать: ``основание больше нуля, значит знак сохраняется''. &
Нужно сравнивать основание с единицей: при \(a>1\) знак сохраняется, при \(0<a<1\) меняется. \\
\addlinespace
Забывать ОДЗ логарифма. &
Всегда сначала требуйте положительность всех аргументов логарифмов. \\
\addlinespace
В дроби включать корень знаменателя в ответ. &
Корни знаменателя всегда выколоты. \\
\addlinespace
Не менять знак при умножении на \(-1\). &
При умножении или делении неравенства на отрицательное число знак разворачивается. \\
\addlinespace
Путать систему и совокупность. &
Система даёт пересечение, совокупность даёт объединение. \\
\addlinespace
Решать дробное неравенство без приведения к нулю. &
Сначала справа должен быть ноль, слева одна дробь. \\
\addlinespace
Забывать, что \(|F(x)|>0\) означает \(F(x)\ne 0\). &
Если модуль стоит под логарифмом, его значение не может быть нулём. \\
\addlinespace
Раскрывать модуль двумя ветвями, не посмотрев на ОДЗ. &
Сначала изучите ограничения: иногда они сразу показывают знак подмодульного выражения. \\
\bottomrule
\end{longtable}

\newpage

\section{Тренировочный блок}

Решите неравенства. Ответ запишите в виде промежутков.

\subsection*{Средний уровень}

\begin{enumerate}
    \item
    \[
        \log_3(x-2)\ge 1.
    \]

    \item
    \[
        \log_2\frac{x+1}{x-3}<1.
    \]

    \item
    \[
        \log_4\frac{x^2-4}{x+3}\ge \frac12.
    \]
\end{enumerate}

\subsection*{Уровень выше среднего}

\begin{enumerate}[resume]
    \item
    \[
        \log_{\frac12}(x^2-4x+3)\le \log_{\frac12}(x-1).
    \]

    \item
    \[
        \log_{\frac34}(x^2-5x+6)>0.
    \]

    \item
    \[
        \log_5\left|\frac{x-1}{x+2}\right|\ge 1.
    \]

    \item
    \[
        \log_2\left|\frac{x+3}{x-1}\right|<2.
    \]

    \item
    \[
        \log_4\frac{4x-5}{|x-1|}>\frac12.
    \]

    \item
    \[
        \log_{\frac13}\left|\frac{x^2-4}{x+1}\right|\le -1.
    \]
\end{enumerate}

\subsection*{Сложный уровень}

\begin{enumerate}[resume]
    \item
    \[
        \log_{\frac12}\left(\frac{|x-2|}{x+1}\right)
        \ge
        \log_{\frac12}(3-x).
    \]
\end{enumerate}

\newpage

\section{Ответы для самопроверки}

\renewcommand{\arraystretch}{1.45}
\begin{table}[h!]
\caption{Ответы к тренировочному блоку}
\centering
\begin{tabularx}{\linewidth}{>{\centering\arraybackslash}p{0.12\linewidth}X}
\toprule
\textbf{№} & \textbf{Ответ} \\
\midrule
1 &
\[
    x\in [5;+\infty)
\]
\\
2 &
\[
    x\in (-\infty;-1)\cup(7;+\infty)
\]
\\
3 &
\[
    x\in \left(-3;\,1-\sqrt{11}\right]\cup
    \left[1+\sqrt{11};+\infty\right)
\]
\\
4 &
\[
    x\in [4;+\infty)
\]
\\
5 &
\[
    x\in \left(\frac{5-\sqrt5}{2};\,2\right)\cup
    \left(3;\,\frac{5+\sqrt5}{2}\right)
\]
\\
6 &
\[
    x\in \left[-\frac{11}{4};-2\right)\cup
    \left(-2;-\frac32\right]
\]
\\
7 &
\[
    x\in (-\infty;-3)\cup
    \left(-3;\frac15\right)\cup
    \left(\frac73;+\infty\right)
\]
\\
8 &
\[
    x\in \left(\frac32;+\infty\right)
\]
\\
9 &
\[
    x\in
    \left(-\infty;\frac{-3-\sqrt{13}}{2}\right]\cup
    \left[\frac{3-\sqrt{37}}{2};-1\right)\cup
    \left(-1;\frac{-3+\sqrt{13}}{2}\right]\cup
    \left[\frac{3+\sqrt{37}}{2};+\infty\right)
\]
\\
10 &
\[
    x\in
    \left[\frac{3-\sqrt{13}}{2};2\right)\cup
    \left(2;\frac{1+\sqrt{21}}{2}\right]
\]
\\
\bottomrule
\end{tabularx}
\end{table}

\end{document}